\section{贪污贿赂罪 - 斡旋受贿、利用影响力受贿罪}
\begin{enumerate}
  \item 斡旋受贿。三方关系，实践中更多。买的中间人的权。原理同普通受贿。
        \begin{enumerate}
          \item \zd{斡旋：说情，给个面子}。
          \item 斡旋者是国家工作人员，不包括单位；为请托人谋取不正当利益；定的罪名还是受贿罪。
          \item 行为方式。
                \begin{enumerate}
                  \item 事前。\al{老板甲 \yjt 财政局长乙 \yjt 税务局长丙 \yjt 非法减免税款}。
                  \item 事后。\al{财政局长乙 \yjt 税务局长丙 \yjt 非法减免甲公司税款 \yjt 乙找老板甲要钱}。
                  \item 间接斡旋。\al{老板甲 \yjt 财政局乙 \yjt 建设局丁 \yjt 税务局丙 \yjt 非法减免税款}。斡旋受贿共犯。
                \end{enumerate}
          \item 和普通受贿区别。\zd{上级不要求是直接领导}。
                \begin{enumerate}
                  \item 上级找下级办事 \yjt 普通受贿。\al{老板甲 \yjt 副县长乙 \yjt 税务局丙 \yjt 非法减免税款}。
                  \item 下级找上级 \yjt 斡旋受贿。\al{老板甲 \yjt 税务局科员乙 \yjt 税务局丙 \yjt 非法减免税款}。
                  \item 同级找同级 \yjt 斡旋受贿。
                \end{enumerate}
          \item 共犯问题。% https://q.uiver.app/?q=WzAsNyxbMCwxLCLogIHmnb/nlLIiXSxbMSwxLCLkuZlcXFxc5Lit6Ze05Lq6Il0sWzIsMSwi5LiZXFxcXOeojuWKoeWxgCJdLFsyLDIsIuWPl+i0v+e9qiJdLFsxLDIsIuaWoeaXi+WPl+i0vyJdLFswLDIsIuWvueacieW9seWTjeWKm+eahOS6uuihjOi0vyBcXGh6IOihjOi0v+e9qiJdLFsxLDAsIui0ouaUv+WxgC/kuJnnmoTlj7jmnLov5oOF5aaHIl0sWzAsMSwiNTBXIl0sWzEsMiwiNDBXIl0sWzIsM10sWzEsNCwi5a6e6KGM54qvIl0sWzAsNV0sWzEsNl0sWzEsMywi5YWx54qvIiwxXV0=
\[
  \begin{tikzcd}[row sep=large]
    & \txt{财政局/丙的司机/情妇} \\
    \txt{老板甲} & \txt{乙\\中间人} & \txt{丙\\税务局} \\
    \txt{对有影响力的人行贿 \hz 行贿} & \txt{斡旋受贿} & \txt{受贿}
    \arrow["50W", from=2-1, to=2-2]
    \arrow["40W", from=2-2, to=2-3]
    \arrow[from=2-3, to=3-3]
    \arrow["\txt{正犯}"{description}, from=2-2, to=3-2]
    \arrow[from=2-1, to=3-1]
    \arrow[from=2-2, to=1-2]
    \arrow["\txt{共犯}"{description}, from=2-2, to=3-3]
  \end{tikzcd}
\]

        \end{enumerate}
  \item 利用影响力受贿。
        \begin{enumerate}
          \item 斡旋者不是国家工作人员，对终端办事的那个官员有影响力的人。
          \item 实质区分：发挥影响力的根据不同。
                \begin{enumerate}
                  \item 私人关系；工作关系/职权关系。可以竞合。
                  \item \zd{买的终端的人的权}，需要终端的人的许诺。
                \end{enumerate}
          \item 共犯关系。同斡旋受贿。
          \item 中间人角色。
                \begin{enumerate}
                  \item 转交。钱不是给中间人的。
                  \item 就是给中间人的。
                  \item 不明说。
                \end{enumerate}
        \end{enumerate}
  \item 对有影响力的人行贿。
  \item 巨额财产来源不明。
\end{enumerate}
